\section{总体分析}
三个问题共享同一条计算主线。问题一首先需要形成可复核的光学评价内核：由日期和时点确定太阳方向与法向直接辐射辐照度，再由太阳入射方向和镜心指向集热器中心的方向确定定日镜姿态，分别计算余弦、大气透射、阴影遮挡、截断和反射效率，最后汇总为镜场热功率与月、年指标。

问题二和问题三不应另行改变评价口径，而应把问题一的计算内核作为候选镜场的统一评价器。问题二的设计变量采用统一镜面尺寸与安装高度；问题三进一步放开逐镜尺寸和高度。两问均需同时处理场界、塔周禁布区、旋转离地、镜间距和$60$ MW额定功率约束。

因此，全文框架按“统一评价内核$\rightarrow$约束可行性$\rightarrow$候选设计比较$\rightarrow$高精度复算”的顺序组织。候选搜索阶段允许使用经记录和验证的近似，但最终表格、Excel提交文件和论文数值必须来自同一套冻结结果。

问题一采用经建模者确认的耦合低差异光线追迹作为正式数值求解器，并以低分辨率因子化光线追迹作为可完成全部输出的基线。主方法和基线共享同一连续光学模型，区别只在积分离散方式；后续问题应沿用问题一已经验证的太阳角、反射、阴影遮挡、有限圆柱截断和功率聚合口径。
